\chapter{Results}

\section{Label Distribution and Agreement}
\label{sec:results-labels}

This section reports the final labels, produced by Claude Sonnet~4.5
(LLM-derived labels), and what they mean for the probe target. Their agreement
with human labels is reported with the labeling method
(Section~\ref{sec:human-audit}): 84 of 100 audited episodes, Cohen's
$\kappa = 0.787$, with no episode labeled Adoption on one side and Retention on
the other.

\begin{table}[htbp]
    \centering
    \caption{Final labels by source (all partitions; LLM-derived labels).
    \emph{Binary}: Adoption and Retention episodes as a share of all episodes.
    \emph{Adoption share}: Adoption as a share of the binary episodes.}
    \label{tab:final-labels}
    \small
    \setlength{\tabcolsep}{4pt}
    \begin{tabular}{lrrrrrrrr}
        \toprule
        Source & Episodes & Adoption & Retention & Agreement & Out & Excluded & Binary & Adoption share \\
        \midrule
        MuSiQue   & 22,132  & 7,481  & 6,403  & 8,068  & 179 & 1  & 62.7\% & 53.9\% \\
        Synthetic & 79,036  & 10,635 & 11,005 & 57,130 & 225 & 41 & 27.4\% & 49.1\% \\
        \midrule
        Total     & 101,168 & 18,116 & 17,408 & 65,198 & 404 & 42 & 35.1\% & 51.0\% \\
        \bottomrule
    \end{tabular}
\end{table}
% Source: results/claude_validation/mixed_final_001/<partition>/raw_results.jsonl

\paragraph{The binary target is balanced.} Of the 35,524 Adoption and
Retention episodes, 51.0\% are Adoption, and the share lies between 50.9\% and
51.6\% in each of the four partitions. Under free-form feedback the same share
was 91.7\% (Section~\ref{sec:problems}), so the changes to the protocol
achieved what they were meant to.

\paragraph{The two sources contribute differently.} Most synthetic episodes
are Agreement (72.3\%), because the Solver and the Critic usually agree on
these questions, as they did in the pilot. The synthetic set is 3.6 times
larger than MuSiQue, so it still supplies 61\% of the binary episodes. MuSiQue
episodes are binary far more often (62.7\%), in line with the Solver's much
lower first-answer accuracy on MuSiQue in the pilot (21\% against 78\%). The
probes therefore learn mostly from synthetic examples, which is why results
are also reported by source.

\paragraph{Out is rare.} Only 404 episodes (0.4\%) are Out, so Claude's weakest
label (precision 60\% in the audit) affects very little data.

\paragraph{Candidate order still matters, in opposite directions by source.}
Each question's order was assigned at random before generation, so within a
source the difference between the two orders measures the effect of order
itself (Table~\ref{tab:order-final}). In MuSiQue episodes, listing the
Critic's answer first raises Adoption by 6.4 points (95\% CI 4.8 to 8.1),
the same direction as in the earlier MuSiQue collection labeled by the
rule-based labels (Section~\ref{sec:problems}). In synthetic episodes it
lowers Adoption by 8.1 points (95\% CI 6.8 to 9.4): the Solver picks the
second-listed candidate more often in both orders (53.2\% when A1 is listed
first, 54.9\% when the Critic is). The two effects differ by 14.5 points
($z = 13.4$), far more than sampling variation would produce. Why they differ
is not established: source was not randomized, and the two sources differ in
hop mix and difficulty. Pooled over both sources, the effects largely cancel
(52.2\% against 49.8\%), which hides them. Order, together with source, is a
signal a probe could use instead of the feedback itself, so every probe and
intervention analysis has to control for both.

\begin{table}[htbp]
    \centering
    \caption{Adoption share of the binary episodes by candidate order, with
    95\% Wilson intervals (final Claude Sonnet~4.5 labels). The order effects
    in the text use normal-approximation intervals for a difference of two
    proportions; the source $\times$ order interaction is the difference
    between the two sources' order effects, with a $z$-test.}
    \label{tab:order-final}
    \begin{tabular}{lll}
        \toprule
        Source & A1 listed first & Critic listed first \\
        \midrule
        MuSiQue   & 50.6\% (49.4--51.8), $n = 6{,}840$  & 57.1\% (55.9--58.2), $n = 7{,}044$ \\
        Synthetic & 53.2\% (52.3--54.1), $n = 10{,}762$ & 45.1\% (44.2--46.1), $n = 10{,}878$ \\
        Both      & 52.2\% (51.5--52.9), $n = 17{,}602$ & 49.8\% (49.1--50.5), $n = 17{,}922$ \\
        \bottomrule
    \end{tabular}
\end{table}
% Source: results/claude_validation/mixed_final_001/<partition>/binary_labels.csv

\section{Objective 1: Pipeline and Activation Dataset}
\label{sec:results-objective1}

The first objective was to build a controlled Solver--Critic interaction
pipeline and an activation dataset for studying it. Both were completed, and
every later experiment uses them.

\paragraph{The pipeline.} One configuration file defines the three roles with
pinned model revisions, the datasets, the layers to record and the decoding
settings. The work is split into chunks that run on separate EC2 machines and
are merged afterwards, with activation row indices rebuilt and checked.
Partitions are fixed by question before any analysis, and related variants of
a question share a group so they always fall in the same partition. Each
stage refuses inconsistent inputs instead of continuing: a label for an
episode that does not exist, an activation index that does not match, or a
retry with different settings.

\paragraph{The dataset.} Table~\ref{tab:dataset} summarizes the final
collection. It holds 101,168 episodes, each with its full text (paragraphs,
A1, the Critic's blind answer and review, A2), correctness against the gold
answer, the candidate order and generation statistics. For each episode the
Solver's activations are stored at nine layers for both prompts, giving
202,336 rows per layer. One in ten episodes is held out for the
interventions, and the rest are split 80/10/10 into training, validation and
test.

\begin{table}[htbp]
    \centering
    \caption{The final collection by partition and source. Binary episodes are
    those labeled Adoption or Retention.}
    \label{tab:dataset}
    \begin{tabular}{lrrrr}
        \toprule
        Partition & Episodes & MuSiQue & Synthetic & Binary episodes \\
        \midrule
        Train        & 72,846  & 15,931 & 56,915 & 25,556 \\
        Validation   & 9,105   & 1,993  & 7,112  & 3,227 \\
        Test         & 9,112   & 1,997  & 7,115  & 3,223 \\
        Intervention & 10,105  & 2,211  & 7,894  & 3,518 \\
        \midrule
        Total        & 101,168 & 22,132 & 79,036 & 35,524 \\
        \bottomrule
    \end{tabular}
\end{table}
% Sources: configs/collection/distributed_final.yaml (partition fractions and seeds);
% results/claude_validation/mixed_final_001/<partition>/raw_results.jsonl (counts).

\paragraph{What the dataset does not contain.} It covers only the natural
feedback condition (Section~\ref{sec:problems}). It stores only the Solver's
activations, and only at the last prompt token
(Section~\ref{sec:activation-caching}); the Critic's internal states are not
recorded. Analyses that need other positions, or the Critic's side of the
exchange, would require a new collection.
\section{Objective 2: SAE Representation}

We compared two SAEs trained directly on the study distribution with a frozen
Gemma Scope SAE and a decoder-adapted version of that pretrained model.  The
custom models used the same 2,560-dimensional Layer~25 input and
12,888-feature dictionary, so their reconstruction and sparsity statistics are
directly comparable.  The pretrained dictionary contained 16,384 features and
was evaluated in the activation units expected by its released checkpoint;
its raw MSE is therefore reported separately rather than compared numerically
with the RMS-normalized custom-model loss.

\subsection{Custom SAE Training}

Both models learned most rapidly in their first few epochs.  The TopK SAE
stopped at epoch~54 under the minimum-improvement and patience rule, with its
lowest validation loss occurring at epoch~51.  The L1 SAE continued improving
through the 200-epoch limit, so its final checkpoint was retained.  Its
training and validation variance explained were 92.3\% and
89.2\%, respectively, a 3.1-percentage-point difference.  The
corresponding TopK values were 92.3\% and 91.9\%, giving nearly identical
training reconstruction but better validation transfer.

The two sparsity mechanisms produced different operating points
(Table~\ref{tab:sae-custom-results}).  On the sealed test split, TopK explained
91.9\% of activation variance with exactly 128 active features per example.
L1 explained 89.1\% while using only 50.8 active features on average.  Thus,
L1 reduced per-example activity by 60.3\% relative to TopK, at a cost of 2.7
percentage points of test variance explained.  This is the central tradeoff
between the custom models: TopK favored reconstruction fidelity, whereas L1
produced the smaller active set.

\begin{table}[htbp]
    \centering
    \caption{Frozen-checkpoint results for the custom SAEs. Reconstruction MSE
    is in RMS-normalized activation units. ``Inactive'' is the fraction of
    dictionary features that never fired on the indicated split; it is not the
    complement of per-example activity.}
    \label{tab:sae-custom-results}
    \begin{tabular}{llrrrr}
        \toprule
        Model & Split & MSE & Variance explained & Mean $L_0$ & Inactive \\
        \midrule
        L1 & Validation & 0.002830 & 89.17\% & 50.64 & 45.82\% \\
        L1 & Test       & 0.002835 & 89.15\% & 50.76 & 45.40\% \\
        TopK & Validation & 0.002119 & 91.89\% & 128.00 & 91.45\% \\
        TopK & Test       & 0.002124 & 91.87\% & 128.00 & 91.35\% \\
        \bottomrule
    \end{tabular}
\end{table}

The dictionary-level activity result qualifies the apparent sparsity
advantage.  Although TopK activates only 128 features for each example, 91.3\%
of its dictionary never fired on the test split.  L1 used fewer features per
example but distributed those activations across a much larger portion of the
dictionary, leaving 45.4\% inactive.  Per-example $L_0$ and dictionary
coverage therefore measure different properties and neither alone establishes
that one representation is more interpretable.

\subsection{Pretrained and Decoder-Adapted SAE}

The frozen pretrained SAE transferred poorly to the project distribution,
explaining 22.5\% of validation variance with an MSE of 6,343.5 in its native
activation units.  Refitting only the reconstruction bias raised variance
explained to 51.2\%, showing that a substantial component of the mismatch was
a distribution-level offset.  The selected epoch-10 decoder-and-bias
checkpoint further reduced validation MSE to 1,689.1 and raised variance
explained to 79.4\% (Table~\ref{tab:sae-pretrained-results}).  This was a 73.4\% MSE
reduction relative to the frozen checkpoint.

\begin{table}[htbp]
    \centering
    \caption{Validation reconstruction for the pretrained SAE variants. MSE is
    in the pretrained model's native activation units. The encoder and its
    feature activations are identical across all three rows.}
    \label{tab:sae-pretrained-results}
    \begin{tabular}{lrrl}
        \toprule
        Variant & MSE & Variance explained & Trainable parameters \\
        \midrule
        Frozen pretrained & 6,343.5 & 22.46\% & None \\
        Bias-only baseline & 3,992.4 & 51.20\% & Reconstruction bias \\
        Decoder adapted & 1,689.1 & 79.35\% & Decoder and bias \\
        \bottomrule
    \end{tabular}
\end{table}

Because adaptation froze the encoder weights, encoder bias, and JumpReLU
thresholds, it did not change which pretrained features fired or the values
provided to the feature-selection probe.  The gain therefore isolates improved
reconstruction of this study's residual stream from improved discovery of
feature identities.  Decoder adaptation made the pretrained representation a
stronger reconstruction baseline and a preserved-basis steering comparator.

Overall, no SAE dominated every criterion.  TopK gave the best custom-model
reconstruction, L1 gave the lowest per-example activity and broader dictionary
coverage, and decoder adaptation preserved pretrained feature identities while
substantially correcting reconstruction mismatch.  These representation-level
results did not define the downstream model set by reconstruction error alone.
TopK was excluded from final downstream tasking; the predictive and causal
experiments used the L1 custom SAE and the pretrained family, including its
adapted variants.
\section{Objective 3: Predictive Features}
\label{sec:res-predictive}

\subsection{Probe Performance}

Every representation supported an above-chance probe of the Solver's
upcoming decision, but none matched the raw activations
(Table~\ref{tab:probe-performance}). The raw layer-25 probe reached 0.843
balanced accuracy. The custom $L_1$ SAE reached 0.773, and the pretrained
family reached about 0.706. The pretrained and decoder-adapted SAEs share an
encoder and therefore give identical probe results; adapting the decoder
changes reconstruction, not which features fire. All SAE probes clearly
exceeded the text-only (0.623), metadata (0.525), and pre-feedback (0.597)
baselines from Section~\ref{sec:layer-selection}. The SAE features therefore
carry decision information that is not visible in the prompt text and that
depends on the feedback, although the sparse bases lose part of what a dense
linear probe can read. This gap is consistent with reports that SAE-based
probes often fall short of probes on raw activations
\cite{kantamneni2025useful}, and the project did not require SAEs to win this
comparison, only that their features carry enough signal to be selected.

\begin{table}[htbp]
    \centering
    \caption{Feature-selection probe on validation. ``Live'' counts latents
    that fire at least once on train. $L_0$ is the mean number of active
    latents per validation example, and FVE is the fraction of variance
    explained by the reconstruction. $^\dagger$Exploratory.}
    \label{tab:probe-performance}
    \begin{tabular}{lrrrr}
        \toprule
        Representation & Live / total & $L_0$ & FVE & Balanced acc.\ [95\% CI] \\
        \midrule
        Raw activations & 2,560 / 2,560 & --- & 1.000 & 0.843 [0.830, 0.856] \\
        Custom $L_1$ SAE & 6,515 / 12,888 & 46.6 & 0.883 & 0.773 [0.758, 0.787] \\
        Pretrained SAE & 3,991 / 16,384 & 112.0 & 0.238 & 0.706 [0.690, 0.722] \\
        Decoder-adapted SAE & 3,991 / 16,384 & 112.0 & 0.770 & 0.706 [0.690, 0.722] \\
        Fully fine-tuned SAE$^\dagger$ & 3,942 / 16,384 & 105.1 & 0.628 & 0.707 [0.691, 0.723] \\
        \bottomrule
    \end{tabular}
\end{table}

\subsection{Offline Ablation Check}

Zeroing the ten targeted features moved the readout's predicted acceptance
far more than zeroing matched random features for every SAE
(Table~\ref{tab:offline-ablation};
Figure~\ref{fig:appendix-feature-selection-proxy}). For the $L_1$ SAE,
targeted ablation lowered predicted acceptance by 40.4 percentage points,
whereas the mean random ablation raised it by 0.1. The decoder-adapted SAE
showed $-21.0$ against $-1.3$. Every SAE's targeted effect had $p\approx10^{-4}$, the
smallest value the permutation test can report. Two results limit how this
should be read. First, the frozen pretrained SAE's reconstruction alone
lowered predicted acceptance from 50.5\% to 25.0\%, reflecting its 0.238
variance explained on this distribution, and its random draws were very
spread out (5th--95th percentile 15.4--37.6\%). Decoder adaptation restored
the unablated rate to 51.0\%. Second, ablating the ten top-ranked raw neurons
had no detectable effect ($-0.17$ percentage points, $p=0.39$). The decision
signal in the raw basis is spread over many neurons, so removing ten of them
changes little, whereas the SAE bases concentrate it enough for a ten-feature
set to matter. This favors the SAE bases for targeted intervention, but the check
is still a probe proxy. Section~\ref{sec:res-causal} tests the same features
in generation.

\begin{table}[htbp]
    \centering
    \caption{Offline ablation on validation: change in probe-predicted
    acceptance (percentage points) relative to each representation's
    unablated reconstruction. Random is the mean over 100 matched draws.
    Observed validation acceptance is 51.0\%. $^\dagger$Exploratory.}
    \label{tab:offline-ablation}
    \begin{tabular}{lrrr}
        \toprule
        Representation & Unablated & Random $\Delta$ & Targeted $\Delta$ [95\% CI] \\
        \midrule
        Raw activations & 50.5\% & $+0.1$ & $-0.2$ [$-0.6$, $+0.2$] \\
        Custom $L_1$ SAE & 50.5\% & $+0.1$ & $-40.4$ [$-41.2$, $-39.5$] \\
        Pretrained SAE & 25.0\% & $-0.3$ & $-9.4$ [$-9.8$, $-9.0$] \\
        Decoder-adapted SAE & 51.0\% & $-1.3$ & $-21.0$ [$-21.4$, $-20.6$] \\
        Fully fine-tuned SAE$^\dagger$ & 50.8\% & $-1.5$ & $-9.4$ [$-9.8$, $-9.1$] \\
        \bottomrule
    \end{tabular}
\end{table}

\subsection{Feature Interpretation}

To describe what the selected features respond to, each was profiled on the
25,556 train episodes in three ways:
\begin{itemize}[nosep]
    \item jointly regressing its activation on twelve episode properties, such
    as the Critic's answer type, candidate order, and the Solver's decision;
    \item projecting its decoder direction through Gemma's unembedding (a logit
    lens);
    \item asking Claude Sonnet~4.5 to describe it from 30 train episodes, then
    scoring that description by its Spearman correlation with the feature's
    activation on 20 validation episodes.
\end{itemize}

The selected features mostly track \emph{what kind of answer the Critic gave},
not deference itself. Across the 80 SAE features, ``the Solver adopted'' was
the strongest property for only four. For eleven of the twenty pretrained
features the strongest property was that the Critic's answer is a date. For
nine of the twenty $L_1$ features it was that the Critic said the question
cannot be answered. The correlation of feature activation with the decision
itself was modest, averaging about $\pm0.1$ (pretrained Adoption and
Retention sets $+0.075$ and $-0.087$; $L_1$ sets $+0.121$ and $-0.061$).
Some individual features read naturally:
\begin{itemize}[nosep]
    \item Pretrained feature 4409 fires on 85\% of episodes, tracks the Critic
    objecting that the question cannot be answered, and promotes tokens such as
    ``option(s)''.
    \item The $L_1$ Retention features 3399 and 8893 track the Solver keeping
    its answer despite such an objection and promote ``unable'' and
    ``insufficient''.
\end{itemize}
Others capture shortcuts. The two top $L_1$ Adoption features, 8677 and 4455,
are most strongly tied to whether the Critic's answer was listed as candidate
A, with opposite signs.

The $L_1$ features were also easier to describe. Their Claude descriptions
achieved a median validation correlation of 0.77, against 0.38 for the
pretrained features and 0.35 for the fully fine-tuned features. As
Section~\ref{sec:res-causal} shows, this ordering is reversed for steering:
the more interpretable $L_1$ features were the weaker causal handles.

\section{Objective 4: Causal Tests}
\label{sec:res-causal}

Of the 1,000 study episodes, 983 were analyzed. Seventeen were excluded
because an edited generation was empty and could not be labeled, and all of
those empty generations came from random-direction edits. The unedited replay
reproduced the stored response and label for all 1,000 episodes. The analyzed
cohort contains 496 episodes whose stored decision was Adoption and 487 whose
stored decision was Retention.

\subsection{Primary Results}

Targeted edits flipped the Solver's decision more often than equal-size
random edits for the two pretrained-family SAEs, and the primary claim holds:
four of the nine prespecified tests survived Holm correction
(Table~\ref{tab:primary-causal}). For the pretrained SAE, the \textsc{both}
arm raised the flip rate from 3.7\% under random edits to 6.2\% ($+2.5$
percentage points, $p_{\mathrm{Holm}}=0.004$). The \textsc{decrease} arm
added 1.7 points ($p_{\mathrm{Holm}}=0.037$). The decoder-adapted SAE
behaved almost identically, with $+2.0$ points for \textsc{both}
($p_{\mathrm{Holm}}=0.007$) and $+1.6$ points for \textsc{decrease}
($p_{\mathrm{Holm}}=0.039$). Its \textsc{increase} arms were positive but did
not survive correction. The custom $L_1$ SAE's \textsc{decrease} arm was the
largest single estimate ($+1.8$ points) but narrowly missed correction
($p_{\mathrm{Holm}}=0.051$). Its \textsc{increase} and \textsc{both} arms were
null.

\begin{table}[htbp]
    \centering
    \caption{Primary causal tests ($n=983$, $r=0.2$). Flip rates are
    percentages of episodes. Random is the mean over three matched draws per
    episode. $\Delta$ is targeted minus random in percentage points, with a
    paired bootstrap 95\% CI. $p$-values are one-sided, Holm-adjusted over the
    nine tests. Bold rows survive at $\alpha=0.05$.}
    \label{tab:primary-causal}
    \begin{tabular}{llrrrr}
        \toprule
        SAE & Arm & Targeted & Random & $\Delta$ [95\% CI] & $p_{\mathrm{Holm}}$ \\
        \midrule
        Custom $L_1$ & decrease & 6.7 & 4.9 & $+1.8$ [$+0.4$, $+3.3$] & 0.051 \\
        Custom $L_1$ & increase & 5.6 & 5.6 & $+0.0$ [$-1.3$, $+1.3$] & 0.999 \\
        Custom $L_1$ & both & 6.4 & 6.5 & $-0.1$ [$-1.7$, $+1.4$] & 0.999 \\
        \textbf{Pretrained} & \textbf{decrease} & 5.1 & 3.4 & $+1.7$ [$+0.4$, $+2.9$] & \textbf{0.037} \\
        Pretrained & increase & 3.9 & 2.7 & $+1.2$ [$+0.2$, $+2.1$] & 0.051 \\
        \textbf{Pretrained} & \textbf{both} & 6.2 & 3.7 & $+2.5$ [$+1.1$, $+3.9$] & \textbf{0.004} \\
        \textbf{Decoder-adapted} & \textbf{decrease} & 4.9 & 3.3 & $+1.6$ [$+0.4$, $+2.9$] & \textbf{0.039} \\
        Decoder-adapted & increase & 4.2 & 3.1 & $+1.1$ [$+0.0$, $+2.3$] & 0.092 \\
        \textbf{Decoder-adapted} & \textbf{both} & 5.6 & 3.6 & $+2.0$ [$+0.7$, $+3.4$] & \textbf{0.007} \\
        \bottomrule
    \end{tabular}
\end{table}

The effects are real but small. Even the strongest arm flips about one
additional episode in forty beyond what an equally large random edit does.
None of the nine prespecified between-SAE contrasts survived correction; the
closest was $L_1$ versus pretrained on the \textsc{both} arm ($-2.6$ points,
$p_{\mathrm{Holm}}=0.059$). No between-arm contrast within an SAE survived
either. One geometric fact explains the null $L_1$ \textsc{both} arm. Its
Adoption and Retention directions are highly aligned (cosine 0.754, against
0.128 for the pretrained SAE), so their difference
$u_{\mathrm{toward}}-u_{\mathrm{away}}$ cancels most of what either set
represents.

The raw-activation control produced larger targeted-minus-random effects
($+4.0$, $+1.7$, and $+4.1$ points for \textsc{decrease}, \textsc{increase}, and
\textsc{both}), but at a cost: its targeted \textsc{decrease} edit left 10.1\%
of answers outside both candidates. In a post hoc comparison, the raw basis
significantly exceeded an SAE only once in nine tests (raw minus $L_1$ on
\textsc{both}, $+4.2$ points, $p_{\mathrm{Holm}}=0.004$). The selected SAE
features steer less strongly than probe-weighted neurons but disrupt
generation much less.

\subsection{Direction, Switches, and Broken Answers}

Effects were stronger when pushing the Solver away from adopting than toward
it (Figure~\ref{fig:appendix-directional-strata}). Averaged over the nine SAE
rows, targeted-minus-random was $+1.7$ points for Adoption-to-Retention and
$+0.9$ for Retention-to-Adoption. Within-representation Holm tests confirmed
this pattern: $L_1$ \textsc{decrease} on Adoption episodes ($+4.0$ points,
$p_{\mathrm{Holm}}=0.006$) and pretrained \textsc{both} on Adoption episodes
($+3.4$, $p_{\mathrm{Holm}}=0.023$) survived, and no SAE row survived in the
Retention-to-Adoption direction. The $L_1$ SAE's near-miss in the pooled
test is thus a one-directional effect diluted by a null in the other
direction.

Flips were usually switches rather than broken answers. Targeted SAE edits
produced few Out answers. The highest were the pretrained and decoder-adapted
\textsc{decrease} arms, at 2.6\% and 3.9\% (against 0.6\% under their random
edits) (Figure~\ref{fig:appendix-out-rate}). Restricting the analysis to
answers that chose one of the two candidates left the SAE effects essentially
unchanged: pretrained \textsc{both} stayed at $+2.5$ points and decoder-adapted
\textsc{decrease} went from $+1.6$ to $+1.8$
(Figure~\ref{fig:appendix-non-out-robustness}). The effects are also not
driven by a few unstable episodes, since the three most-flipped episodes
account for only 5.5\% of targeted flips
(Figure~\ref{fig:appendix-flip-distribution}). They are concentrated where the
Solver was least certain: most episodes had a baseline first-token
probability of at least 0.999, and the flips and the targeted-random gap came
mainly from the less confident remainder
(Figure~\ref{fig:appendix-confidence-sensitivity}).

Two examples show what a switch looks like. In an Adoption-to-Retention
example, for the question ``In what year was the architect of the Marnhill
Civic Theatre born?'', the Solver first answered 1927, the Critic proposed
1884, and the stored second answer adopted 1884. The pretrained
\textsc{decrease} edit produced 1927, while all three random edits produced
1884. In a Retention-to-Adoption example, ``When did the performer of
\emph{Out from Under} release her first album?'', the Solver kept 1999 against
the Critic's 2000. The decoder-adapted \textsc{both} edit produced 2000, and
all random edits kept 1999. In both cases the edited answer is a fluent choice
of the other candidate, not degraded text.

\subsection{Reconstruction Control}

Passing the activation through an SAE without editing it changed some
decisions. Reconstruction-only flip rates were 1.6--2.2\% for Adoption
episodes and 1.6--2.3\% for Retention episodes across the four SAEs, against
0\% for the unedited replay. The custom $L_1$ SAE had the smallest relative
reconstruction error (0.049) but not the fewest flips, so flip rates did not
simply track reconstruction error. These rates set a noise floor for how
sensitive the decision is to small perturbations at this position. They do
not contaminate the edit effects, because edits were applied to the original
activation rather than to the reconstruction, and every comparison is
against random edits of the same size.

\subsection{Exploratory Comparisons}

The fully fine-tuned SAE did not differ reliably from the pretrained SAE
(differences of $-1.3$ to $+0.8$ points, none significant;
Figure~\ref{fig:appendix-full-finetuning}). The LLM-edited pretrained SAE
gave the largest targeted-minus-random effects of any representation
($+8.5$ to $+9.6$ points) and beat its mechanical-rotation control by 2.2 to
3.0 points (Holm $p\le0.005$). The comparison with Out rates qualifies this
(Figure~\ref{fig:appendix-llm-edit-flips-vs-out}). The LLM edit's targeted Out
rate was 7--12\% against 2.5--7.5\% for the mechanical edit, and most extra Out
answers were ``The evidence is insufficient.'' The LLM-designed directions
push generation harder rather than aiming it more precisely. Finally, in a
post hoc split by data source, three tests survived Holm on MuSiQue episodes
($n=371$) and none on synthetic episodes ($n=612$). The causal effects
therefore rest mainly on the natural questions.

\section{Earlier Exploratory Runs: Layer 33 and Layer 25 Steering}
\label{sec:res-exploratory}

The final protocol is the fourth causal design we ran. The earlier runs used
an older collection, rule-based behavior labels, and custom SAEs trained
before layer selection had identified layer~25. Their numbers are not
comparable to the final results, and we report them because each one
determined a design choice.

\paragraph{Probe-scored suppression (October 6 and 14).} The first pipeline
zeroed or amplified ($z\mapsto3z+1$) individual SAE latents and read the
outcome from a probe applied to the re-encoded latents rather than from
generation. These runs most likely used the layer-33 custom SAE that was the
initial single-layer benchmark, although their configuration does not record
the layer. The ``flip rate'' under reconstruction alone (0.63 and 0.23 in the
two runs) was as large as under targeted suppression or amplification
(0.61--0.63 and 0.23--0.27) and as under random features. The apparent
effects came entirely from passing through the SAE. \emph{Lesson:} every
causal comparison needs a reconstruction control and a random control, and
outcomes should be measured in generation.

\paragraph{Zero-ablation in generation (October 18).} We then zeroed the
targeted features at layers 25 and 33 while the Solver generated its answer,
on 517 episodes whose baseline replay matched exactly. At layer~25, $L_1$
targeted ablation reduced acceptance by 29.6 points (95\% CI 25.3--33.7).
However, the edit removed 93\% of the activation norm, against 0.8\% for
random ablations, and the outputs were broken rather than switched. Of 458
episodes that originally adopted, 83 produced a third answer, 70 were
unresolved, and only one actually retained the Solver's first answer. 6.96\%
of all answers were empty, and others degenerated into strings such as
``ikorododododo''. At layer~33, no SAE showed a targeted effect. \emph{Lesson:}
zeroing a dominant feature set is a large, uncontrolled perturbation. Zero
ablation was kept only as the offline probe check of
Section~\ref{sec:feature-selection}, and broken outputs had to be reported
separately from decision changes.

\paragraph{Clamped steering (October 20).} The next design clamped each
targeted feature to a multiple $s\in\{-2,+2,+5\}$ of its mean active value,
with seven frequency-matched random draws per condition. At layer~25 the $L_1$
SAE appeared to steer powerfully, $+73.0$ points targeted-minus-random (95\%
CI 68.8--76.8), with $+25.3$ at layer~33, while the pretrained SAE showed
nothing. An audit traced the large effects to edit size. Targeted edits were
8--37 times larger in norm than random ones: for $L_1$ at layer~25, $2.83$
versus $0.08\lVert x\rVert$ at $s=-2$. One always-on feature supplied 86\%
(layer~25) or 91\% (layer~33) of the set's activation, negative clamps went
far outside the observed range, and answers almost never changed for edits
below about $0.4\lVert x\rVert$. \emph{Lesson:} targeted and random edits must
have equal size, and feature magnitude must not determine edit size. This led
to unit-direction edits with a fixed $r\lVert x\rVert$ and a larger pool of
random draws.

\paragraph{Equal-size edits (November 3).} With edit sizes matched, a TopK
SAE run at layer~25 showed that turning all ten selected features off still
reduced acceptance ($-5.1$ points, $p_{\mathrm{Holm}}<0.001$) but required an
edit of 86\% of $\lVert x\rVert$. A matched 1\% edit moved the probe by $-2.5$
points but left generated answers unchanged ($+0.01$). This run introduced
the explicit split of flips into \emph{switches} (a change to the other
candidate) and \emph{broken} answers (a change to Out). It showed that the
all-off edit produced a real switch toward Retention together with broken
answers, while removing the single dominant feature produced mostly broken
answers. \emph{Lesson:} the probe is far more sensitive than generation, so
the edit size must be chosen from generated answers on a held-out pilot, and
Out must be reported alongside flips. Both became part of the final protocol
(Sections~\ref{sec:causal-methods} and~\ref{sec:metrics}).
\section{Overall Answer to the Research Question}
